% GENERATED FILE. Edit canonical lessons, then run npm run generate:books.
% canonical-source-hash: fnv1a64:c3c0a75b048fe004
% canonical-lessons: SA-C76-rain, SA-C76-cold, SA-C76-heat, SA-C76-spring, SA-C76-summer

\chapter{Rain, Cold and Heat}
\label{ch:sa-rain-cold-heat}
\hlchaptermodality{76}

\begin{chapteropening}
\textbf{By the end of this chapter:} \emph{I can name rain, the cold, the heat, spring and summer.}

Complete the last lesson of chapter 76: I can name rain, the cold, the heat, spring and summer.
\end{chapteropening}

\section[\texorpdfstring{\sk{वर्षा} (varṣā)}{वर्षा (varṣā)}]{\sk{वर्षा} (varṣā) — rain}

\label{lesson:SA-C76-rain}



\begin{quote}
Before the new one: say the Sanskrit for an aeroplane, then the Sanskrit for he goes in.
\end{quote}

\subsection*{You'll want to know: \sk{वर्षा}}
\textbf{\sk{वर्षा}} (\emph{varṣā}) — \textquotedblleft{}rain\textquotedblright{}, and in the plural the rainy season.

\begin{cousinweb}
\emph{Varṣā} is built on the root √vṛṣ, \textquotedblleft{}to rain\textquotedblright{} (MW, CDIAL 11396). Its sister \emph{varṣa}, \textquotedblleft{}rain; a year\textquotedblright{}, counted in rainy seasons, lives on as Hindi \emph{baras}, \textquotedblleft{}year\textquotedblright{} (CDIAL 11392).
\end{cousinweb}

\begin{grammarlens}[title={What you've built}]
Water from the sky.
\end{grammarlens}

\subsection*{Your turn}
\begin{itemize}
\raggedright
  \item \emph{Say it:} \emph{varṣā}
  \item \emph{Say it:} \emph{varṣā}, once more
  \item \emph{Say it:} say \emph{varṣā}
  \item \emph{Recall:} say the Sanskrit for an aeroplane, then the Sanskrit for he goes in, then say \emph{varṣā} again
\end{itemize}

\begin{tcolorbox}[breakable,colback=teal!4,colframe=teal!35!black,title={Before you move on}]
What does \sk{वर्षा} mean? (\textquotedblleft{}Rain\textquotedblright{}.) Say it once more.
\end{tcolorbox}
\section[\texorpdfstring{\sk{शीतम्} (śītam)}{शीतम् (śītam)}]{\sk{शीतम्} (śītam) — cold}

\label{lesson:SA-C76-cold}



\begin{quote}
Before the new one: say the Sanskrit for he goes in, then the Sanskrit for rain.
\end{quote}

\subsection*{You'll want to know: \sk{शीतम्}}
\textbf{\sk{शीतम्}} (\emph{śītam}) — \textquotedblleft{}cold\textquotedblright{}. \textbf{\sk{शीतम्} \sk{अस्ति}} — \textquotedblleft{}it is cold\textquotedblright{}.

\begin{cousinweb}
\emph{Śīta} is built on the root √śyai, \textquotedblleft{}to congeal, freeze\textquotedblright{} (MW, CDIAL 12485), and is already in the Rigveda.
\end{cousinweb}

\begin{grammarlens}[title={What you've built}]
The weather of the snow.
\end{grammarlens}

\subsection*{Your turn}
\begin{itemize}
\raggedright
  \item \emph{Say it:} \emph{śītam}
  \item \emph{Say it:} \emph{śītam}, once more
  \item \emph{Say it:} say \emph{śītam asti}
  \item \emph{Recall:} say the Sanskrit for he goes in, then the Sanskrit for rain, then say \emph{śītam} again
\end{itemize}

\begin{tcolorbox}[breakable,colback=teal!4,colframe=teal!35!black,title={Before you move on}]
What does \sk{शीतम्} mean? (\textquotedblleft{}Cold\textquotedblright{}.) Say it once more.
\end{tcolorbox}
\section[\texorpdfstring{\sk{उष्णम्} (uṣṇam)}{उष्णम् (uṣṇam)}]{\sk{उष्णम्} (uṣṇam) — heat}

\label{lesson:SA-C76-heat}



\begin{quote}
Before the new one: say the Sanskrit for rain, then the Sanskrit for the cold.
\end{quote}

\subsection*{You'll want to know: \sk{उष्णम्}}
\textbf{\sk{उष्णम्}} (\emph{uṣṇam}) — \textquotedblleft{}heat\textquotedblright{}, \textquotedblleft{}hot\textquotedblright{}. \textbf{\sk{उष्णम्} \sk{अस्ति}} — \textquotedblleft{}it is hot\textquotedblright{}.

\begin{cousinweb}
\emph{Uṣṇa} is a Rigvedic word for \textquotedblleft{}hot\textquotedblright{}, which MW files under the root √uṣ, \textquotedblleft{}to burn\textquotedblright{} (CDIAL 2389).
\end{cousinweb}

\begin{grammarlens}[title={What you've built}]
The other side of cold.
\end{grammarlens}

\subsection*{Your turn}
\begin{itemize}
\raggedright
  \item \emph{Say it:} \emph{uṣṇam}
  \item \emph{Say it:} \emph{uṣṇam}, once more
  \item \emph{Say it:} say \emph{uṣṇam asti}
  \item \emph{Recall:} say the Sanskrit for rain, then the Sanskrit for the cold, then say \emph{uṣṇam} again
\end{itemize}

\begin{tcolorbox}[breakable,colback=teal!4,colframe=teal!35!black,title={Before you move on}]
What does \sk{उष्णम्} mean? (\textquotedblleft{}Heat\textquotedblright{}.) Say it once more.
\end{tcolorbox}
\section[\texorpdfstring{\sk{वसन्तः} (vasantaḥ)}{वसन्तः (vasantaḥ)}]{\sk{वसन्तः} (vasantaḥ) — spring}

\label{lesson:SA-C76-spring}



\begin{quote}
Before the new one: say the Sanskrit for the cold, then the Sanskrit for the heat.
\end{quote}

\subsection*{You'll want to know: \sk{वसन्तः}}
\textbf{\sk{वसन्तः}} (\emph{vasantaḥ}) — \textquotedblleft{}spring\textquotedblright{}, the first \textbf{\sk{ऋतुः}} of the year, when the \textbf{\sk{पुष्पम्}} opens.

\begin{cousinweb}
MW explains \emph{vasanta} as \textquotedblleft{}the brilliant season\textquotedblright{}, from the root √vas, \textquotedblleft{}to shine\textquotedblright{} (CDIAL 11439). It is inherited: MW compares Greek \emph{ear} and Latin \emph{vēr}, \textquotedblleft{}spring\textquotedblright{}, the source of English \textbf{vernal}.
\end{cousinweb}

\begin{grammarlens}[title={What you've built}]
The flower season.
\end{grammarlens}

\subsection*{Your turn}
\begin{itemize}
\raggedright
  \item \emph{Say it:} \emph{vasantaḥ}
  \item \emph{Say it:} \emph{vasantaḥ}, once more
  \item \emph{Say it:} say \emph{vasantaḥ}
  \item \emph{Recall:} say the Sanskrit for the cold, then the Sanskrit for the heat, then say \emph{vasantaḥ} again
\end{itemize}

\begin{tcolorbox}[breakable,colback=teal!4,colframe=teal!35!black,title={Before you move on}]
What does \sk{वसन्तः} mean? (\textquotedblleft{}Spring\textquotedblright{}.) Say it once more.
\end{tcolorbox}
\section[\texorpdfstring{\sk{ग्रीष्मः} (grīṣmaḥ)}{ग्रीष्मः (grīṣmaḥ)}]{\sk{ग्रीष्मः} (grīṣmaḥ) — summer}

\label{lesson:SA-C76-summer}



\begin{quote}
Before the new one: say the Sanskrit for the heat, then the Sanskrit for spring.
\end{quote}

\subsection*{You'll want to know: \sk{ग्रीष्मः}}
\textbf{\sk{ग्रीष्मः}} (\emph{grīṣmaḥ}) — \textquotedblleft{}summer\textquotedblright{}, the \textbf{\sk{ऋतुः}} of \textbf{\sk{उष्णम्}}.

\begin{cousinweb}
\emph{Grīṣma} is already in the Rigveda. The native grammarians derive it from the root √gras, \textquotedblleft{}to devour\textquotedblright{} (MW). MW also compares an Irish word for fire, a comparison CDIAL 4391 does not take up, so its origin is not settled.
\end{cousinweb}

\begin{grammarlens}[title={What you've built}]
The hot season.
\end{grammarlens}

\subsection*{Your turn}
\begin{itemize}
\raggedright
  \item \emph{Say it:} \emph{grīṣmaḥ}
  \item \emph{Say it:} \emph{grīṣmaḥ}, once more
  \item \emph{Say it:} say \emph{grīṣmaḥ}
  \item \emph{Recall:} say the Sanskrit for the heat, then the Sanskrit for spring, then say \emph{grīṣmaḥ} again
\end{itemize}

\begin{tcolorbox}[breakable,colback=teal!4,colframe=teal!35!black,title={Before you move on}]
What does \sk{ग्रीष्मः} mean? (\textquotedblleft{}Summer\textquotedblright{}.) Say it once more.
\end{tcolorbox}
